\documentclass[11pt]{article}
\usepackage[utf8]{inputenc}
\usepackage[margin=1in]{geometry}
\usepackage{amsmath, amssymb, amsthm}
\usepackage{enumitem}
\usepackage{algorithm}
\usepackage{algpseudocode}
\usepackage{fancyhdr}

\setlength{\headheight}{14pt}
\pagestyle{fancy}
\fancyhf{}
\lhead{CS 521: Algorithms}
\chead{Homework 4}
\rhead{Razz}
\cfoot{\thepage}

\newcounter{problem}
\newcommand{\problem}[1]{\stepcounter{problem}\section*{Problem \theproblem: #1}}
\newcommand{\bigO}{\mathcal{O}}

\begin{document}

\begin{center}
  {\Large \textbf{Homework 4: Dynamic Programming and Greedy Algorithms}} \\[4pt]
  Due Friday, 30 October 2026
\end{center}

\problem{Longest increasing subsequence}

Given an array $a_1, \dots, a_n$ of integers, give an $\bigO(n \log n)$ algorithm that returns the length of a longest strictly increasing subsequence.

\paragraph{Solution.} Maintain an array $t$ in which $t[\ell]$ is the smallest value that ends an increasing subsequence of length $\ell$ among the elements seen so far. The array $t$ is strictly increasing, so each new element can be placed by binary search.

\begin{algorithm}[h]
  \caption{Longest increasing subsequence}
  \begin{algorithmic}[1]
    \Function{LIS}{$a_1, \dots, a_n$}
      \State $t \gets$ empty array
      \For{$i \gets 1$ \textbf{to} $n$}
        \State $\ell \gets$ the least index with $t[\ell] \ge a_i$, found by binary search
        \If{no such index exists}
          \State append $a_i$ to $t$
        \Else
          \State $t[\ell] \gets a_i$
        \EndIf
      \EndFor
      \State \Return the length of $t$
    \EndFunction
  \end{algorithmic}
\end{algorithm}

\emph{Correctness.} By induction on $i$, after processing $a_i$ the entry $t[\ell]$ is the smallest possible last element of an increasing subsequence of length $\ell$ in $a_1, \dots, a_i$. \emph{Running time.} Each iteration performs one binary search on an array of length at most $n$, so the total is $\bigO(n \log n)$.

\problem{Knapsack with repetition}

There are $n$ item types with integer weights $w_i > 0$ and values $v_i$, and a knapsack of capacity $W$. Each type may be used any number of times.

\paragraph{Solution.} Let $K(c)$ be the best value achievable with capacity $c$. Then $K(0) = 0$ and
\[
  K(c) = \max\Bigl( K(c - 1), \ \max_{i \,:\, w_i \le c} \bigl( K(c - w_i) + v_i \bigr) \Bigr) \qquad \text{for } c \ge 1 .
\]
Filling the table for $c = 1, \dots, W$ takes $\bigO(nW)$ time and $\bigO(W)$ space. This is pseudo-polynomial: $W$ is exponential in the length of its encoding.

\problem{Interval scheduling}

Prove that choosing, at each step, the compatible interval that finishes first yields a maximum set of pairwise disjoint intervals.

\begin{proof}
Let $g_1, \dots, g_k$ be the greedy choice and $o_1, \dots, o_m$ an optimal solution, both sorted by finishing time. We show by induction that $f(g_j) \le f(o_j)$ for every $j \le k$. The case $j = 1$ holds because greedy picks the interval that finishes first overall. If $f(g_{j-1}) \le f(o_{j-1})$, then $o_j$ starts after $g_{j-1}$ finishes, so $o_j$ was available to greedy, which therefore chose an interval finishing no later.

If $m > k$, the interval $o_{k+1}$ starts after $f(o_k) \ge f(g_k)$ and would have been picked by greedy, a contradiction. Hence $m = k$.
\end{proof}

\problem{Huffman codes}

For the frequencies $a\colon 45$, $b\colon 13$, $c\colon 12$, $d\colon 16$, $e\colon 9$, $f\colon 5$, a Huffman code assigns
\begin{center}
  \begin{tabular}{c|cccccc}
    Symbol & $a$ & $b$ & $c$ & $d$ & $e$ & $f$ \\ \hline
    Codeword & 0 & 101 & 100 & 111 & 1101 & 1100
  \end{tabular}
\end{center}
with expected length $(45 \cdot 1 + 13 \cdot 3 + 12 \cdot 3 + 16 \cdot 3 + 9 \cdot 4 + 5 \cdot 4) / 100 = 2.24$ bits per symbol.

\end{document}
